\section{What did the evidence establish?}\label{sec:experiments}
Two prospective studies put the same question to executable services: when the mathematics says an observation is sufficient, does a real invocation carry the promised information, authority, price and work effect? The earlier studies isolate each ingredient and expose cases in which a simple competitor wins. Some use native policy engines or signed artifacts. Others deliberately construct small services with known ground truth. The original manuscripts and the governing mathematical revisions are identified in the \href{https://github.com/Perceptive-Closure/PC-X/blob/664f90e2d2e2c941bed36e43240811cfc33ef0ac/PAPER/SOURCE-MANIFEST.md}{paper source manifest}. The observations below are attributable to their archived workloads. The earlier Perceptive Closure manuscript and the seventeen PC-X packages share methods and artifacts, so pooling them as independent evidence for a population claim would misstate the design. Appendix~\ref{app:experiments} gives denominators, comparator scopes and replay instructions.

\paragraph{When a successful run hides its mechanism.}
In the original manuscript, a projection of Apache Polaris left two different resource explanations for a repair. A separate specimen using one unchanged AuthzForce policy and engine gave different decisions on two world-specific requests. Admitting the missing attribute closed the certificate at cost one. Freezing evidence made this repair impossible. The same manuscript included 135 constructed PC-TAU tasks. Only 3,309 of 5,436 scheduled language-model conditions resolved. Transport failures accounted for 1,982 and budget stops for 145. Frequencies among resolved conditions therefore describe a selected subset. In X01 and X02, giving the controller hidden state understated costs, while discarding correlations made some identified coordinates appear ambiguous. X06 and X07 tested authorization and certificate records. On X07 the native signed-field baseline tied the proposed procedure on all eight verdicts.

\paragraph{Being selective can cost more.}
In X03 and X04, buying a minimum certificate or adapting later reads to earlier reports sometimes reduced the worst price. A five-site routing case needed two adaptive units and six static units. X05 represented \(2^{192}\) contracts compactly, yet both target and identity acquisition cost 677 on that service. Lazy search took longer than direct partitioning on 20 of 28 jointly solved cases, eager search on 22. On X12, the direct full-acquisition comparator resolved ten of twelve units for 786,444 configured operations; the anchored method resolved six for 1,048,804. The examples establish possible savings and identify mechanisms that can eliminate them.

\paragraph{Checking the measured event.}
X08--X11 separated benchmark observations from actual outcomes, tested evaluation-aware behavior and introduced monitor or intervention changes. X09's greedy comparator tied, and both unsupported scoring-shift cases were falsely certified. X11's disturbed probe preserved public transcripts while increasing actual risk by over 66 percentage points in both contexts. X13's isolated-component comparison falsely certified three of four composed ports. X14 explicitly required useful work: the governed arm was safe on all ten tickets and completed eight, whereas permanent refusal was safe on ten and completed none. X15 independently replayed X01--X14. Its custody audit recorded same-second verification and completion for X08 and reliance on the parent ledger for the X11 opening time.

\paragraph{Can doing the work supply the next observation?}
The X16 study \citep{expx16} installed a source whose lawful routes follow paths in a correlated graph. Each actual job wrote durable work records. Across eight profiles the preregistration predicted the target verdict for 178 admissible installed paths and 70 numerical or infeasible outcomes for the comparator policies and solvers. All matched their forecasts. There were 22 primary held-out installations, ten additional held-out control installations, and 649 native trials by the independent role. One profile had zero paid cost when the first genuine job disclosed the needed route. A fixed schedule failed there. Suppressing that receipt made the common adaptive service impossible even though both known installations could complete with work budget four. With the installed path supplied, a native comparator read the actual guard rows and executed the least-work known-mode schedule with zero paid audit charge. The common policy faced the harder task of choosing before it knew the path. Separate explicit, full-score and compressed solvers compared the same unknown-path requirement. On another profile extending the move horizon from three to four reduced the adaptive price from three to two. A matched explicit solver was faster by median runtime on six of the eight profiles. These are results for the signed finite source and its checked projections.

\paragraph{Can the same run buy information, withstand a disclosed verdict and finish safely?}
The X17 study \citep{expx17} used two separate runners and decoders on one native service. A signed public randomness round selected eight held-out rational settings after source, forecast and custody commitments. Each role ran five policies on each setting, with 4,096 attempts per policy and setting. All 327,680 main attempts produced valid records. A separate exact optimizer checked the matched lower bound and its attainment. The primary and independent mean per-run PC-X probe charges were 0.583797 and 0.582767, respectively, against \(7/12\). High or skip cost \(2/3\) in the exact restricted menu. Low or skip and no paid probe failed the uniform error target. A more informative direct observation achieved \(1/6\) on its richer menu. The guard cost \(1/4\) per installation separately.
\figpc{risk}{0.96}{X17 selective violation risk among actual issued runs. Points show observed proportions. Bars give the registered simultaneous intervals for each world and role. Every upper endpoint is below \(3/8\). Data come from the two separately decoded native runs, with positive issuance observed in every cell.}{fig:risk}
The denominator is actual issuance, including positive decisions in truly negative worlds. Confidence concerns repeated sampling by these runners. Deployment risk concerns physical violations on issued attempts. All sixteen PC-X world/run intervals supported the \(3/8\) limit. The tighter \(5/16\) limit was supported in six of eight worlds per role. The \(5/32\) positive-only limit was supported in all four positive worlds per role. The uniform label-error threshold was supported in four of eight per role, and the positive issuance and utility floors in two of four per role. A separately registered analysis used a signed score, the observed violation-and-issuance indicator minus \(3/8\) times the issuance indicator. Its interval supported the \(3/8\) selective-risk requirement in six of eight primary and five of eight independent settings. It asks about the same physical cap through a different statistical calculation. The all-world issuance floor received support in six of eight per role. These are threshold decisions at the registered simultaneous confidence level. All 640 point forecasts lay inside their corresponding intervals. At the separate \(2/25\) equivalence tolerance, 314 of 320 primary and 312 of 320 independent comparisons established equivalence; fourteen stayed inconclusive. All 320 same-cell interval pairs overlapped.

Both roles used the same host and service, with separately committed seeds, programs and decoders. Each role compared 1,024 nominal/actual diagnostic pairs. All six disagreements per role matched the physical surcharge event on the common prefix. All nineteen registered control effects reproduced on their specified witnesses. An invalid post-verdict control exposed a scoring difference between success on one leg and utility over the whole run. The exact source result assumes the good guard, declared event decoder, issued-set denominator and retained prelude. At prelude length six a grid certificate took about 269 milliseconds with counted states and 1,191 with explicit histories on the measured host. At length one the counted representation used more peak traced memory. At length 256 it constructed 33,153 prelude states per route; there was no timing comparison with the full tree at that length. Independently checked grid intervals narrowed from \([7/12,121/192]\) to \([7/12,233/384]\). Their upper endpoints are attainable cost bounds; the separate analytic proof gives the exact optimum. The samples are repeated attempts in eight settings on one constructed service. These qualifications matter for any extrapolation to a different service.
